
%!TEX root = ../template.tex
%%%%%%%%%%%%%%%%%%%%%%%%%%%%%%%%%%%%%%%%%%%%%%%%%%%%%%%%%%%%%%%%%%%
%% 6 - Results.tex
%% NOVA thesis document file
%%
%% Chapter with results
%%%%%%%%%%%%%%%%%%%%%%%%%%%%%%%%%%%%%%%%%%%%%%%%%%%%%%%%%%%%%%%%%%%
% chktex-file 24
\typeout{NT FILE 6 - Results}%

\chapter{Results}
\label{cha:results}

This chapter presents the results of the performance evaluation of the 4G and 5G networks implemented in the testbed described in Chapter~\ref{cha:implementation}. 
% ============================================================================
% BASE STRUCTURE -- proposed storyline, drafted from Test_Data/analysis/FINDINGS.md,
% FINDINGS_RERUN_2026-08-02.md and FINDINGS_RERUN_2026-08-03.md. Fill in numbers from
% those files / the fig*.pdf figures they reference. Each section below notes which
% round(s) the data comes from -- be explicit about this in the prose, since the 5G
% config (TDD pattern, rx_gain) differs between rounds and the numbers aren't all
% directly comparable to each other.
%
% Framing note for the whole chapter: there is no single "best" round to headline as
% "the" 5G result. Round 2 (rx_gain 70) and Round 3 (rx_gain 60) each win at one distance
% and lose badly at the other -- that's not a flaw in the presentation, it's the
% headline finding of Section 5 below. Consider saying this explicitly near the start of
% the chapter so the reader doesn't wonder why later sections mix numbers from different
% rounds.

% --- 1. Throughput overview (headline comparison) ---------------------------------
% The chapter's opening result: 4G baseline vs 5G, TCP/UDP x UL/DL, at each location.
% Source: FINDINGS.md S1 (Round 1, fig1_throughput.pdf) for the 4G baseline and the
% "before" picture; FINDINGS_RERUN_2026-08-02.md S1 (fig7_tdd_rebalance.pdf) and
% FINDINGS_RERUN_2026-08-03.md S1 (fig10_rxgain_tradeoff.pdf) for the two corrected
% configurations. Recommend presenting Round1 -> Round2 -> Round3 as a progression per
% location/direction rather than one static table, since the story IS the progression.
% Headline numbers to hit:
%   - Round 1: 5G uplink up to 21x worse than 4G at 10m (worst point).
%   - Round 2: TDD rebalance alone recovers 5G uplink to ~match 4G at 2m and 10m, but
%     downlink regresses as the airtime trade-off (quantify: DL 0.45-0.79x of Round-1 5G).
%   - Round 3: rx_gain 60 pushes 2m uplink to 2x the 4G baseline, but 10m collapses to
%     0.17x of Round 2's value. State plainly: neither Round 2 nor Round 3 is "the" 5G
%     result, they bracket the same problem from both sides (see Section 5).
%   - Location B (6m) throughput: excluded from Round 2 (unusable, see Section 4);
%     recovered in Round 3 -- only report Round 1 and Round 3 numbers for this location.

\section{General comparison between 4G and 5G}
\label{sec:4g-vs-5g}

This section presents a comparison of the 4G baseline against the 5G network across the three test rounds described in Section~\ref{sec:test-methodology}, focusing on the evolution of performance metrics. Because the 5G configuration changed between rounds, namely the TDD structure and the \texttt{rx\_gain} setting, these results are presented as a progression rather than a single snapshot, tracking how each change affected throughput at each location and in each direction. As will become clear over the following sections, no single round can be presented as the definitive 5G result: the configuration that performs best at short range is not the same as the one that performs best at long range, a trade-off examined in more detail later in this chapter.

In Figures \ref{fig:tp-tcp-dl} through \ref{fig:tp-udp-ul}, the mean throughput results for TCP and UDP in both downlink and uplink directions are presented for each test location, with black lines representing the range between Min/Max throughput values between each repetition for that combination. These figures mainly illustrate the performance differences between 4G and 5G across the three rounds of testing, with the results showing that the 5G network can outperform the 4G baseline in certain configurations and locations, but also underperform in others.

\begin{figure}
    \centering
    \includegraphics[width=\textwidth]{TCP_DL.pdf}
    \caption{TCP downlink throughput at each test location.}
    \label{fig:tp-tcp-dl}
\end{figure}

\begin{figure}
    \centering
    \includegraphics[width=\textwidth]{UDP_DL.pdf}
    \caption{UDP downlink throughput at each test location.}
    \label{fig:tp-udp-dl}
\end{figure}

\begin{figure}
    \centering
    \includegraphics[width=\textwidth]{TCP_UL.pdf}
    \caption{TCP uplink throughput at each test location.}
    \label{fig:tp-tcp-ul}
\end{figure}

\begin{figure}
    \centering
    \includegraphics[width=\textwidth]{UDP_UL.pdf}
    \caption{UDP uplink throughput at each test location.}
    \label{fig:tp-udp-ul}
\end{figure}


% --- 2. Uplink airtime and TDD decomposition ---------------------------------------
% THE ARGUMENT: the 3:1 pattern was a hard structural cap on uplink, independent of signal quality. Rebalancing to 2:2 lifted the cap and paid for it in downlink airtime. This is the cleanest quantitative result in the dataset, and -- unlike the raw throughput numbers -- it is measured on both sides and survives the SNR confound below.
%
% LEAD WITH THE DECOMPOSITION, NOT THE FIGURE (FINDINGS.md S2). Separate the two factors by counting grants rather than Mb/s:
%   - Factor A, airtime. 3:1 gives one uplink slot per 2.5 ms = a 400/s ceiling; the measured rate is 372-394 grants/s. 4G FDD gets one uplink subframe per ms, ~975/s. That is a 2.5x structural penalty no amount of signal quality can recover.
%   - Factor B, link budget. Bits per grant. At 2 m the 5G and 4G payloads are effectively identical (13 634 vs 13 429 bits); at 10 m the 5G payload collapses 8.7x.
%   Multiply them: 2.5 x 1.0 = 2.6x deficit at 2 m (measured 2.59x), 2.5 x 8.7 = 21.5x at 10 m (measured 21.5x). The decomposition is exact at both distances. This is what earns the section its own heading.
%
% THE REBALANCE DELIVERED EXACTLY WHAT WAS PREDICTED (FINDINGS_RERUN_2026-08-02.md). 2:2 gives two uplink slots per 2.5 ms = 800/s, and the scheduler grants the full 800/s at both 2 m and 10 m. The x2 airtime factor is therefore measured, not inferred. Airtime also predicts downlink x0.75 (4 -> 3 DL slots); observed TCP DL is x0.56 at A and x0.44 at C, i.e. downlink lost somewhat more than airtime alone accounts for.
%
% WHAT FIGURE~\ref{fig:tdd-structure} SHOWS. Only one variable moves per step: TDD between the first two points, rx_gain between the last two. Say that explicitly, or a reader in a section titled "TDD decomposition" will attribute the second step to airtime as well.
%   DL/UL ratio, TCP:   A: 5.84 -> 1.32 -> 0.89     C: 17.55 -> 0.39 -> 2.48
%   At A the two directions converge monotonically and finish nearly balanced (22.4 DL vs 25.2 UL) -- the best balance anywhere in the dataset. At C they do not: rx_gain 60 gives back most of the uplink Round 2 recovered (11.60 -> 2.77 Mb/s, a 4.2x regression) and the ratio swings past 1 to 2.48. So write "close at short range", not "close". The Location C behaviour is not a blemish on the result -- it is the rx_gain trade-off that Section~\ref{sec:uplink-quality-rxgain} builds to, and this figure is its clearest evidence. Forward-reference it rather than explaining it here.
%
% CONFOUND THAT HAS TO BE STATED (FINDINGS_RERUN_2026-08-02.md). Between Rounds 1 and 2, at identical gain settings, uplink SNR rose +5.4 dB at A and +22.4 dB at C -- documented and still unexplained. So the x19.9 uplink gain at C is mostly NOT airtime, which alone predicts x2.0. Do not present the Round 1 -> 2 throughput change as a clean measurement of the TDD effect. The grants/s counting above is the part that survives, because grants/s is independent of SNR: the confound lives entirely in bits-per-grant.
%
% OPTIONAL SECOND FIGURE -- the measured slot map (data_v2/tdd_airtime.csv; plot-ready in data_v2/fig_tdd_tradeoff.xlsx, sheet SlotMap). PUSCH appears only in slot 4 under 3:1 and in slots 3-4 under 2:2: the structural ceiling straight from the scheduler's own counters, and the one piece of evidence the confound cannot touch. Three traps:
%   (a) use the slot POSITIONS, not the PRB magnitudes -- the 3:1 unit was an uplink-only run while the 2:2 unit aggregates every test at that location, so the PRB totals (58.5 vs 42.1) are not a capacity ratio and must not be quoted as one;
%   (b) for the same reason the 3:1 PDSCH column is near-zero from absent traffic, not from unavailable slots;
%   (c) Round 1 has no airtime data at all -- its logs carry no Scheduler-cell METRICS lines -- which is why the slot map uses the Round 3 pair (same location, same gain, same session) rather than the Round 1 vs Round 2 pair used in the figure.
%
% ALSO STATE:
%   - Location B is absent from the figure: Round 2 at B is the unusable metal-stand data.
%   - The Round 3 TDD isolation run (R3_5G_C_TDD31) cannot settle this. It is TCP uplink only, so it has no downlink leg and cannot show a trade-off at all, and it ran at rx_gain 60, which desensitises Location C. Report it as inconclusive rather than as a contradiction, and flag "redo at rx_gain 70" as future work.
%   - The first configuration is n=1 per condition (Round 1); the other two are n=3.
\section{TDD rebalance effects on uplink airtime and throughput}
\label{sec:tdd-rebalance}

% REMINDER (framing): the throughline for this section, and for the chapter, should not be
% "5G beats/loses to 4G" -- it's "the uplink deficit is diagnosable and largely a
% configuration choice, not a fundamental 5G limitation." The 3:1 TDD pattern was inherited
% from an old worked example and never revisited (see Implementation ch.); that's a real,
% fixable misconfiguration, and the Round 2 rebalance recovering most of the uplink is the
% legitimate positive result to lean on here -- not a claim that 5G outperforms 4G.
% Don't spend effort trying to spin the numbers themselves positively; the diagnostic value
% of "we found the cause and fixing it worked" is the stronger, more defensible thesis
% contribution than a scoreboard result.

One of the troubles identified on the first round of testing was that the 5G uplink throughput was capped by the TDD structure configuration implemented in the testbed, which, as mentioned in Section~\ref{sec:test-methodology}, allocated TDD timeslots disproportionately in favor of the downlink. The figures in the section show that with this configuration, the 5G downlink reigned with 68.8\% of the total airtime, compared with 22.9\% for the uplink.

We can verify this asymmetry by examining the raw metrics from the gNB trace files, checking the scheduler's actual uplink grant rate against the theoretical maximum allowed by the TDD configuration.

On the radio configuration file, shown in Listing~\ref{lst:cell_cfg}, the \texttt{cell\_cfg} block sets the \gls{SCS} to 30\,kHz. As explained in Chapter~\ref{cha:5g-fundamentals}, at this \gls{SCS}, a slot is 0.5\,ms long, and the 3:1 TDD pattern sets a 5-slot transmission period with one uplink slot, giving a total period of 2.5\,ms. This means that the maximum number of uplink slots that can be scheduled per second is $1000\,\text{ms} / 2.5\,\text{ms} = 400$ uplink slots, regardless of the radio link quality.

\begin{lstlisting}[language=yaml,backgroundcolor=\color{gray!20},frame=single,framerule=0.8pt,rulecolor=\color{black},basicstyle=\ttfamily\small,caption={\texttt{gnb.yaml} TDD configuration for the 3:1 pattern used in Round 1.},label={lst:cell_cfg}]
cell_cfg:
  dl_arfcn: 627340
  band: 78
  channel_bandwidth_MHz: 20
  common_scs: 30
  plmn: "00101"
  tac: 7
  pci: 1
\end{lstlisting}

Checking the raw metrics from an uplink test extracted from one of the gNB trace files in Round 1 (excerpt shown in Listing~\ref{lst:ul-grant-ceiling}), we see that the \texttt{ok} and \texttt{nok} columns sum to 400 transmissions for each second, confirming that the scheduler was granting the full 400 uplink slots/s allowed by the 3:1 TDD pattern. This confirms that the uplink throughput was indeed capped by the TDD configuration, rather than by the radio link quality.

\begin{lstlisting}[backgroundcolor=\color{gray!20},frame=single,framerule=0.8pt,rulecolor=\color{black},basicstyle=\ttfamily\small,columns=fixed,keepspaces=true,caption={Console trace extracted from the \texttt{gnb.log} file from the Location A (2\,m) TCP uplink test in Round~1. The \texttt{ok} and \texttt{nok} columns are the counts of successful and failed uplink transmissions in that one-second reporting period; \texttt{pusch}/\texttt{rsrp} are the reported SINR and RSRP, \texttt{brate} the resulting bit rate.},label={lst:ul-grant-ceiling}]
|-------------------------------UL-----------------------------
| pusch  rsrp  ri  mcs  brate   ok  nok  (%)    bsr     ta  phr
|  18.3 -17.0   1   14  7.77M  382   18   4%   700k   240n    3
|  19.1 -16.2   1   15  7.68M  376   24   6%   700k   227n    0
|  19.7 -15.6   1   16  6.25M  308   92  23%   700k   234n   -3


\end{lstlisting}

The \texttt{ok} and \texttt{nok} columns sum to $382+18=400$ transmissions for the first second, as does the seconds after. Averaged over every one-second reporting row of the full Round~1 uplink runs whose grant count reaches at least 90\% of the technology's structural ceiling --- a rule applied uniformly below to exclude call-setup and TCP-slow-start rows without hand-picking any individual second --- this settles to 400.0 grants/s at Location A and 399.7 grants/s at Location C, matching the same maximum established before. The 4G baseline, measured the same way from the \texttt{enb.log} traces, grants 1000.0 and 992.5 subframes/s at the same two locations. This equates to a 2.5$\times$ reduction ($1000.0/400.0$ at Location A, $992.5/399.7$ at Location C) no amount of signal quality could recover, since it is imposed before a single symbol is transmitted.

A second factor --- the payload carried per grant --- also contributes to the deficit, but it tracks distance-dependent link quality (MCS and link adaptation) rather than the TDD pattern itself, so it is examined together with that cause in Section~\ref{sec:uplink-quality-rxgain} instead of here. A working comparison of 4G and 5G payload per grant, developed while investigating this, is kept at the end of this file pending that rewrite.

% REMINDER: this is the section's actual good-news result -- lean on it. The 3:1 -> 2:2
% rebalance recovering most of the uplink demonstrates the Factor A deficit above was a
% fixable config choice, not a hardware or protocol limit. Don't undercut it by also trying
% to spin the payload-per-grant point positively (see the staged comparison at the end of
% this file) or by folding the rx_gain/OVL problems from Section~\ref{sec:uplink-quality-rxgain}
% into "misconfiguration" too -- those look more like open-source RAN stack immaturity
% (harder to tune than a mature 4G stack) than a single fixable knob, and deserve that more
% honest framing rather than "my config was wrong."
On the subsequent rounds of testing, the TDD pattern was rebalanced to a symmetrical 2:2 pattern, which allowed two full downlink and uplink slots, with the flexible slot also adjusted to be symmetrical, per period. This change was expected to improve uplink throughput, at the cost of downlink throughput, making their performance almost equal, which Figure~\ref{fig:tdd-structure} presents. Only one configuration variable changes between each pair of rounds in the figure (the TDD pattern between rounds 1 and 2, and \texttt{rx\_gain} between rounds 2 and 3) so the two steps should be read as evidence for two separate effects rather than a single continuous trend.

The pattern used from Round~2 onward, shown in Listing~\ref{lst:tdd-22-cfg} of Chapter~\ref{cha:implementation}, doubles \texttt{nof\_ul\_slots} to 2 within the same 5-slot, 2.5\,ms period, raising the structural ceiling to $2\times400=800$ uplink transmissions per second. Round~2 onward records srsRAN's native timestamped scheduler metrics rather than the periodic console dump used in Round~1, so the raw evidence takes a different shape: Listing~\ref{lst:ul-800-grant-ceiling} shows the relevant uplink fields from three consecutive one-second \texttt{gnb.log} entries for the Location A TCP uplink test (the full lines also report downlink and per-HARQ statistics, omitted here for space).

\begin{lstlisting}[backgroundcolor=\color{gray!20},frame=single,framerule=0.8pt,rulecolor=\color{black},basicstyle=\ttfamily\small,columns=fixed,keepspaces=true,caption={Uplink fields extracted from three consecutive \texttt{gnb.log} scheduler-metrics entries for the Location A (2\,m) TCP uplink test in Round~2 (2:2 pattern).},label={lst:ul-800-grant-ceiling}]
15:14:46  ul_mcs=14  ul_brate=14.1Mbps  ul_nof_ok=692  ul_nof_nok=108
15:14:47  ul_mcs=17  ul_brate=16.3Mbps  ul_nof_ok=647  ul_nof_nok=153
15:14:48  ul_mcs=19  ul_brate=18.1Mbps  ul_nof_ok=638  ul_nof_nok=162
\end{lstlisting}

\texttt{ul\_nof\_ok} and \texttt{ul\_nof\_nok} sum to exactly 800 in every row shown, and the same holds throughout the saturated portion of the run at both Location A and Location C, confirming the scheduler reaches the full doubled ceiling regardless of distance --- this is what the grant rate alone predicts, measured directly rather than inferred from throughput.

\begin{figure}[htbp]
    \centering
    \includegraphics[width=\textwidth]{TDD_structure_comparison.pdf}
    \caption{TCP throughput at Locations A (2\,m) and C (10\,m) across the three 5G
    rounds. The x-axis shows the specific configuration used in each round.}
    \label{fig:tdd-structure}
\end{figure}

We see on the figure that initially, as expected, when using the 3:1 TDD pattern, the uplink throughput on both location A and C was significantly lower than the downlink throughput. On round 2, after rebalancing the TDD pattern to 2:2, the uplink performance rose to nearly match the downlink in location A, surpassing the downlink throughput at location C. We also note that the downlink throughput decreased in both locations, as expected.

The same airtime accounting also predicts the downlink share should fall to three-quarters of its Round 1 value ($\times0.75$); the observed downlink throughput fell further still, to $\times0.56$ at Location A and $\times0.44$ at Location C, so some factor beyond airtime alone also affected the downlink.

Finally, on round 3, after adjusting the \texttt{rx\_gain} to 60: at location A (2 m), the overall network throughput increased significantly, as this change decreased the number of \texttt{ovl} events on the uplink, thus stabilizing the network's performance. However, at location C, further away from the radio, the uplink throughput decreased significantly as well, due to the lower \texttt{rx\_gain} desensitizing the receiver and reducing the link budget. This change was expected as a trade-off between the two locations, and it highlights the challenge of finding a single configuration that optimizes performance across all distances.

We also see that in this final configuration, the throughput between both directions is closest to being equal at short range: at Location A the two directions converge to within a narrow margin (22.4 vs.\ 25.2\,Mb/s), the best balance observed anywhere in the dataset, while at Location C the \texttt{rx\_gain} reduction gives back most of the uplink gain Round 2 had recovered, swinging the downlink-to-uplink ratio past parity to 2.48. This Location C behaviour is not a contradiction of the TDD result but the \texttt{rx\_gain} trade-off that Section~\ref{sec:uplink-quality-rxgain} examines in detail.
% --- 3. Uplink quality vs distance: SINR, PHR, and the rx_gain trade-off -----------
% Source: FINDINGS.md S3 (Round 1 PHR/SINR, and the corrected-vs-superseded discussion of
% the ~33dB deficit -- keep the correction framing, don't restate the disproven power-cap
% explanation as fact) + FINDINGS_RERUN_2026-08-02.md S3 (front-end compression / OVL
% evidence at rx_gain 70) + FINDINGS_RERUN_2026-08-03.md S1 and S3 (the rx_gain 60 sweep
% results, and the PHR "anomaly" resolved as a PRB-allocation artifact -- worth a
% sentence so a reader doesn't think there's still a mystery). This section should end on
% the fixed-gain-vs-AGC finding, quoted near-verbatim from FINDINGS_RERUN_2026-08-03.md:
% "rx_gain 70 is correct at 10m and overloads at 2m; rx_gain 60 is correct at 2m and
% desensitises at 10m -- no single value works." This is probably the single most
% citable sentence in the whole results chapter.

% --- 4. Reliability: RLFs, re-attaches, and the Location B obstruction -------------
% Source: FINDINGS.md S5 (Round 1: 9 RLFs at 2m from missed HARQ-ACKs, 0 in 4G) +
% FINDINGS_RERUN_2026-08-02.md S2 (Round 2: Location B unusable, 18 re-attaches -- tie
% this to the metal-stand obstruction already documented in Section~\ref{sec:test-round-2})
% + FINDINGS_RERUN_2026-08-03.md S0/S5 (Round 3: attach behaviour "perfect" at every
% location, 0 RLFs, stand removed). This section can double as evidence that the stand
% hypothesis from the Methodology chapter is plausible, since its removal correlates
% with the churn disappearing.

% --- 5. TCP/uplink coupling: why downlink didn't benefit from 5G -------------------
% Source: FINDINGS.md S4. One of the more interesting findings -- uplink BLER (53.7% at
% 10m TCP DL vs 3.7% at the same distance/UDP) throttles the TCP downlink via lost ACKs,
% even though the DL radio itself can sustain far more (proven by the UDP DL result at
% the same distance/config). FINDINGS_RERUN_2026-08-03.md S5 has a nice confirming
% coda: fixing the uplink in Round 3 also lifted 2m downlink TCP (16.5 -> 22.4 Mb/s).

% --- 6. Unloaded latency ------------------------------------------------------------
% Source: FINDINGS_RERUN_2026-08-03.md S6 (fig11_unloaded_latency.pdf), Round 3 only --
% state plainly this was not measured in Rounds 1-2 and there is no 4G ping baseline to
% compare against. Headline: 15.7-16.0ms floor, distance-independent (as expected --
% propagation delay over 10m is negligible). The sawtooth/phase-beat analysis (RTT
% cycling with the ~19ms uplink scheduling opportunity) is a nice, non-obvious result if
% there's room for it, but the floor number alone is the one a reader will expect to see.

% --- 7. Bufferbloat: the cost hiding inside the downlink numbers -------------------
% Source: FINDINGS_RERUN_2026-08-02.md S8 (discovery, Round 2) + FINDINGS_RERUN_2026-08-03.md
% S6 "Idle versus loaded" (Round 3, quantified against Section 6's floor: 35-81x RTT
% inflation once a TCP download starts, from a 6.17MB RLC queue). Recommend placing this
% right after Section 6 -- it reframes every downlink throughput number reported earlier
% in the chapter ("10 Mb/s downlink" vs "10 Mb/s downlink at 2.7s RTT" are different
% claims) and works well as a late "twist" rather than a footnote. State clearly this was
% diagnosed but not fixed (RLC resize was deprioritised out of both lab sessions) --
% frame as a concrete, well-evidenced future-work item, not an unresolved failure.

% --- 8. Control-plane latency -------------------------------------------------------
% Source: FINDINGS.md S8 (Round 1 only, fig6_control_plane.pdf). 5G SA attach ~2.2x
% slower end-to-end than 4G (166ms vs 360ms RRC-request-to-bearer), but the air-interface
% stages are comparable or faster -- the gap is procedural (more NAS round trips, split
% registration/PDU-session-setup) and the SBI-based core shows up directly in a 5x
% slower NG setup (9.5ms vs 1.9ms). Not re-measured in Rounds 2-3; say so.

% --- 9. Software RAN feasibility: processing time vs slot deadline -----------------
% Source: FINDINGS.md S7 (Round 1, fig4_processing_time.pdf). 5G p99 processing time
% occasionally exceeds the 500us slot budget (up to 1.4x over) with real logged
% consequences (late DL data, RF underflows); 4G never comes close to its 1ms budget.
% Caveat honestly: different host machines, so state both specs and don't oversell it as
% a clean apples-to-apples CPU comparison. FINDINGS_RERUN_2026-08-02.md S5 notes
% real-time health improved after the TDD rebalance (less DL duty cycle, less PHY load)
% -- worth a sentence tying back to Section 2/3.

% --- 10. Operational lessons: logging cost --------------------------------------
% Source: FINDINGS.md S9 (Round 1). 2.81GB / 1341s at info level, ~7.5GB/hour
% extrapolated to continuous operation -- ties directly to the "leave a working testbed
% for future projects" objective from the Introduction/Implementation chapters, so it's
% worth keeping even though it's not a performance result. Concrete recommendation:
% warning level / hex_max_size 0 for normal operation.

% --- 11. Smaller findings (one paragraph or a summary table) -----------------------
% Candidates, pick what's not already covered incidentally above: 256QAM usage on both
% stacks (FINDINGS.md S10), UL HARQ using rv=0 only / no incremental redundancy
% (FINDINGS.md S6 -- ties into Section 3's BLER discussion, consider merging there
% instead), downlink PRB occupancy saturated in both systems, 4G RF-overflow-at-startup
% as a benign artefact, TCP-vs-UDP protocol sensitivity being 5G-specific.

% --- Open items to state honestly, not paper over -----------------------------------
% - The original ~33dB power deficit's real cause (Section 3) is still unconfirmed --
%   don't let Section 1's later fixes imply it was ever explained, only that later
%   config changes (TDD, gain) improved on it empirically.
% - The 22dB unexplained SINR jump at 10m between the July and August campaigns
%   (FINDINGS_RERUN_2026-08-02.md S4) -- explicitly flagged there as "probably not a
%   config effect." Don't claim the TDD change alone explains the 10m uplink recovery.
%   Worth a sentence in Section 1 or 2 as a limitation.
%   TODO once the 03-Aug data is copied from the flash drive per SESSION_LOG_2026-08-03.md
%   sec:test-round-3 header: check whether it settles this.

% =====================================================================================
% STAGED FOR REFACTOR -- move into Section~\ref{sec:uplink-quality-rxgain} (pulled out of
% sec:tdd-rebalance, which now only covers the grants/s (Factor A) story). This is a
% working comparison of 4G vs.\ 5G payload per grant at Locations A and C, all numbers
% verified against Test_Data/analysis/scripts/ul_grant_decomposition.py (90%-of-ceiling
% saturation rule, stated inline below -- keep that rule if this moves, it's what makes
% the numbers reproducible). Key results if rewriting from scratch instead of moving this
% verbatim: 5G payload/grant at Location A averages 11,282 bits vs. 13,429--13,482 for 4G
% (a 1.2x gap, NOT near-parity -- 5G's channel here is 20MHz vs. 4G's 10MHz, so even a
% 1.2x gap is a real per-Hz shortfall, don't spin it positively); at Location C it's
% 1,485 vs. 12,931 (8.7x). The Location-A listing also shows the same run dropping to the
% same MCS-0 floor Location C sits at permanently -- useful evidence that Location A isn't
% immune to the link-adaptation problem, just less exposed to it. Also carries the SNR
% confound note (Round 1 -> Round 2 raw throughput comparison is contaminated by an
% unexplained SINR rise, 5.4dB at A / 22.4dB at C, so don't present it as a clean TDD
% measurement) which arguably belongs in sec:uplink-quality-rxgain too since it's about
% interpreting throughput through a link-quality lens, not the TDD pattern.
% =====================================================================================

% Another factor to note is the payload carried per grant. Dividing the successfully delivered bit rate from the \texttt{brate} column by the number of successful grants in a given second gives the mean payload size per grant for that second. Listing~\ref{lst:ul-4g-payload} shows an excerpt from the 4G \texttt{enb.log} trace at Location A, where each row grants the full 1000 uplink subframes per second, as the FDD frame structure allows, and payload size works out to roughly 15\,000, 14\,000 and 13\,000 bits per grant.

% \begin{lstlisting}[backgroundcolor=\color{gray!20},frame=single,framerule=0.8pt,rulecolor=\color{black},basicstyle=\ttfamily\small,columns=fixed,keepspaces=true,caption={Console trace excerpt from the 4G \texttt{enb.log} file for the Location A (2\,m) TCP uplink test in Round~1.},label={lst:ul-4g-payload}]
% | pusch  pucch  phr  mcs  brate   ok  nok  (%)    bsr
% |  16.4    n/a   18   17    15M 1000    0   0%   135k
% |  14.7    n/a   19   16    14M 1000    0   0%   141k
% |  13.2    n/a   19   15    13M 1000    0   0%   138k
% \end{lstlisting}

% From the equivalent 5G \texttt{gnb.log} trace at Location A, shown in Listing~\ref{lst:ul-5g-2m-payload}, the same columns give payloads of roughly 11\,200, 11\,600, 8\,000, 2\,500 and 1\,300 bits per grant across five consecutive seconds.

% \begin{lstlisting}[backgroundcolor=\color{gray!20},frame=single,framerule=0.8pt,rulecolor=\color{black},basicstyle=\ttfamily\small,columns=fixed,keepspaces=true,caption={Console trace excerpt from the 5G \texttt{gnb.log} file for the Location A (2\,m) TCP uplink test in Round~1, five consecutive seconds.},label={lst:ul-5g-2m-payload}]
% | pusch  rsrp  ri  mcs  brate   ok  nok  (%)    bsr     ta  phr
% |  13.8 -21.1   1    7  4.48M  400    0   0%   700k   262n   -7
% |  13.8 -21.8   1    8  4.59M  397    3   0%   700k   260n   -3
% |  10.8 -24.3   1    5  3.17M  396    4   1%   700k   234n   -2
% |   1.8 -29.6   1    1   993k  391    9   2%   700k   211n    0
% |  -8.1 -33.5   1    0   535k  400    0   0%   700k   251n   -3
% \end{lstlisting}

% Those five rows already swing by 8.4$\times$ from first to last, which makes clear that a handful of rows cannot stand in for a run average, as MCS is not constant even while every second stays saturated, so the payload per grant moves with it. Within those five seconds, MCS collapses from 7 down to 0 and \texttt{brate}/\texttt{ok} falls from 4.48\,M/400 to 535\,k/400 --- the same MCS-0, roughly-535k-bit-rate floor Listing~\ref{lst:ul-5g-10m-payload} shows as the steady state at Location C. Location A is not immune to the link-adaptation problem examined in Section~\ref{sec:uplink-quality-rxgain}; it just spends most of its time above that floor instead of pinned to it. Applying a 90\%-of-ceiling saturation rule -- keeping only one-second reporting rows whose grant count reaches at least 90\% of the technology's structural ceiling, to exclude call-setup and TCP-slow-start rows without hand-picking any individual second -- and dividing the total delivered bits by the total successful grants across every kept second of the run, the 5G payload per grant at Location A averages 11\,282 bits, against 13\,482 bits for 4G at the same location --- a 1.2$\times$ difference, not the near-parity a couple of high-throughput rows might suggest, though still close enough that Location A is link-budget-limited far less than Location C. This should not be read as ``5G nearly matches 4G per grant'': the 5G channel in this testbed is 20\,MHz vs.\ 4G's 10\,MHz, so double the spectrum for only a 1.2$\times$-closer payload is a per-Hz shortfall, not parity, once a reader does that math themselves. This trace format does not report the PRB allocation per grant, so that cannot be confirmed as the cause either way.

% At Location C (10\,m), the picture is far starker: Listing~\ref{lst:ul-5g-10m-payload} shows the 5G trace at that distance still granting the full ceiling every second (\texttt{ok}\,=\,400), but at MCS~0 --- the lowest possible modulation and coding scheme --- with \texttt{brate}/\texttt{ok} down to roughly 1\,300--1\,340 bits per grant, compared to the steady-state 4G payload at the same distance (12\,931 bits per grant, essentially unchanged from Location A). This tracks the MCS directly: the grant count is unchanged from Location A (still 400/s, the same Factor A ceiling), but MCS falls from single digits to 0, the lowest possible modulation order and code rate, so far fewer bits fit in each grant regardless of how many resource blocks it spans. This trace format does not report the number of resource blocks per grant, so the PRB allocation itself cannot be ruled out as a contributing factor, but the MCS collapse alone is enough to explain a payload reduction of this size. Averaged over the full run with the same saturation rule, the 5G payload per grant collapses to 1\,485 bits, an 8.7$\times$ reduction, while the 4G payload is unchanged.

% \begin{lstlisting}[backgroundcolor=\color{gray!20},frame=single,framerule=0.8pt,rulecolor=\color{black},basicstyle=\ttfamily\small,columns=fixed,keepspaces=true,caption={Console trace excerpt from the \texttt{gnb.log} file for the Location C (10\,m) TCP uplink test in Round~1 (5G). The scheduler still grants the full 400 uplink slots/s, but at the lowest MCS, with a far smaller payload per grant than at Location A.},label={lst:ul-5g-10m-payload}]
% | pusch  rsrp  ri  mcs  brate   ok  nok  (%)    bsr     ta  phr
% |  -2.7 -36.0   1    0   535k  400    0   0%   700k   334n  -13
% |  -6.3 -35.9   1    0   536k  400    0   0%   700k   332n  -11
% |  -7.4 -36.3   1    0   535k  400    0   0%   700k   285n  -12
% \end{lstlisting}

% Multiplying the two factors gives a predicted throughput deficit of $2.5\times1.2=3.0\times$ at Location A and $2.5\times8.7=21.6\times$ at Location C. Measured directly from the same runs, the actual throughput deficit is $3.1\times$ at Location A and $21.6\times$ at Location C: an exact match at Location C, and a close one at Location A, where the small residual gap comes from grant rate and payload per grant being two independently averaged quantities (a mean over seconds and a mean over grants) that do not multiply back to the mean throughput ratio with perfect precision when both vary together within a run. It is this decomposition, not the raw throughput gap, that justifies treating the TDD pattern as a distinct and independently correctable factor from the link-budget problems examined in Section~\ref{sec:uplink-quality-rxgain}.

% One caveat applies to comparing Round 1 and Round 2 by throughput directly: at identical gain settings, uplink SINR rose by 5.4\,dB at Location A and by 22.4\,dB at Location C between the two rounds, for reasons that remain unexplained. The raw uplink throughput gain between rounds therefore cannot be attributed to the TDD change alone --- the $19.9\times$ increase measured at Location C, for instance, far exceeds the $\times2$ that airtime alone predicts. The grants-per-second accounting above is unaffected by this confound, since the achieved grant rate does not depend on the SINR reported for each grant; the confound lives entirely in the bits-per-grant term.

